% Mục 4.3. Nguồn: include/counter.hpp. Lệnh máy: c++ -O2 -S (Apple clang 21), mặc định (LSE),
% -Xclang -target-feature -Xclang -lse, và -arch x86_64.
Counter bắt đầu từ 0. \texttt{increment()} cộng 1 và trả về giá trị \emph{trước} khi tăng. $M$ lần tăng
phải trả về đúng $0, \dots, M - 1$, mỗi giá trị một lần, và theo thứ tự thời gian thực
(linearizable~\cite{herlihy1990linearizability}). Ba cài đặt là \PCname{counter_lock<L>} (biến
\texttt{long} trong critical section của mỗi lock ở phần 1), \PCname{counter_cas} (lock-free, có biến thể
backoff) và \PCname{counter_faa} (wait-free). Trên \PCname{apple-m5} có \PCval{apple-m5}{ncounters} counter.
\begin{lstlisting}
long old = value_.load(relaxed);                   // a guess; the CAS checks it
while (!value_.compare_exchange_weak(old, old + 1, Order, relaxed)) {
    ++failed_;                                     // thread_local, failure path only
    backoff.pause();
}
return old;
// fetch-and-add: return value_.fetch_add(1, Order);
\end{lstlisting}

Thứ tự mặc định là \texttt{relaxed}, vì tính atomic của phép đọc-sửa-ghi và thứ tự toàn phần của mọi phép
ghi lên một biến atomic đúng với mọi memory order. Counter cũng không bảo vệ dữ liệu nào khác. Biến thể
\PCname{faa_seqcst}, \PCname{cas_seqcst} đo cái giá của \texttt{seq\_cst}. Counter dùng lock không có
atomic riêng. Nó đúng nhờ cặp release--acquire của lock.

Apple clang bật LSE (lệnh atomic đơn của ARMv8.1) mặc định. \texttt{fetch\_add} là một lệnh
\texttt{ldadd} (\texttt{ldaddal} với \texttt{seq\_cst}), không vòng lặp. CAS là một lệnh \texttt{cas}
(\texttt{casal}), chỉ hỏng khi giá trị thật sự đã đổi, nên mỗi lần \texttt{cas\_fail} là một xung đột thật.
Khi tắt LSE, cả hai thành \texttt{ldxr}/\texttt{stxr} (LL/SC), và CAS weak có thể hỏng giả. Trên x86-64
hai thứ tự sinh cùng một lệnh (\texttt{lock xadd}, \texttt{lock cmpxchg}).

Giá trị atomic nằm riêng một cache line. Trong \PCname{counter_lock}, lock và giá trị nằm trên hai line
để thread chờ không làm mất hiệu lực line của giá trị. Vì vậy mỗi lần tăng ghi lên hai line. Số lần CAS hỏng
được đếm trong biến \texttt{thread\_local} để không thêm một line bị tranh chấp. Về giới hạn, counter CAS
chỉ lock-free nên một thread có thể bị đói. \PCname{counter_lock} mang theo giới hạn của lock bên trong.
